\section{Proofs}\label{app:proofs}

% A numbered remark environment local to the appendices; defined only once.
\ifcsname appremark\endcsname\else
  \theoremstyle{remark}
  \newtheorem{appremark}{Remark}[section]
  \theoremstyle{plain}
\fi

This appendix proves the results stated in the main text and adds five remarks,
each with its proof, that the main text uses informally. Coordinates are real
numbers, sizes and orientations are fixed, and touching rectangles are legal.
Hard macro $i$ has center $(x_i,y_i)$, width $w_i$ and height $h_i$ on a
$L_x\times L_y$ canvas; for a pair $i\ne j$ write
$\bar w_{ij}=(w_i+w_j)/2$ and $\bar h_{ij}=(h_i+h_j)/2$.

\subsection{Feasible set and coordinate slices}

\begin{proof}[Proof of Proposition~\ref{prop:geometry}]
Let $\mathcal B$ be the set of placements that satisfy every canvas-containment
inequality $w_i/2\le x_i\le L_x-w_i/2$, $h_i/2\le y_i\le L_y-h_i/2$ and every
fixed-coordinate equality. For an unordered pair $e=\{i,j\}$ of hard macros, the
pair does not overlap exactly when one of four closed halfspaces contains the
placement:
\begin{equation*}
 H_{e,1}=\{x_j-x_i\ge\bar w_{ij}\},\quad
 H_{e,2}=\{x_i-x_j\ge\bar w_{ij}\},\quad
 H_{e,3}=\{y_j-y_i\ge\bar h_{ij}\},\quad
 H_{e,4}=\{y_i-y_j\ge\bar h_{ij}\}.
\end{equation*}
With $\mathcal E$ the set of pairs and $m=|\mathcal E|$, the feasible set is
$\mathcal F=\mathcal B\cap\bigcap_{e\in\mathcal E}\bigcup_{b=1}^4H_{e,b}$. A finite
intersection of finite unions equals the union, over all ways of choosing one member
from each union, of the intersections of the chosen members. Choosing branch
$\sigma(e)$ for every pair gives
\begin{equation*}
 \mathcal F=\bigcup_{\sigma\in\{1,2,3,4\}^{\mathcal E}}P_\sigma,\qquad
 P_\sigma=\mathcal B\cap\bigcap_{e\in\mathcal E}H_{e,\sigma(e)}.
\end{equation*}
Each $P_\sigma$ is the intersection of finitely many closed halfspaces and
hyperplanes, hence a closed convex polyhedron. The containment inequalities bound
every coordinate, so $P_\sigma$ is bounded and therefore compact. There are $4^m$
choices of $\sigma$; different choices may give overlapping, identical or empty
polyhedra, so $4^m$ is an upper bound on the number of distinct pieces.
\end{proof}

\begin{appremark}[Legal coordinate slices]\label{rem:slice}
Fix a legal placement and vary only $x_i$ for a movable hard macro $i$. Let
\begin{equation*}
 S_i=\bigl\{j\ne i\text{ hard}:\ |y_i-y_j|<\bar h_{ij}\bigr\}
\end{equation*}
be the macros that overlap $i$'s horizontal band. The legal values of $x_i$ form
\begin{equation}
 I_i=\Bigl[\tfrac{w_i}{2},\,L_x-\tfrac{w_i}{2}\Bigr]\setminus
     \bigcup_{j\in S_i}\bigl(x_j-\bar w_{ij},\,x_j+\bar w_{ij}\bigr),
 \label{eq:app-slice}
\end{equation}
a finite union of closed intervals, possibly including single points.
\end{appremark}
\begin{proof}
Only the pairs that contain $i$ change. For $j\notin S_i$, the placement already
satisfies $H_{e,3}$ or $H_{e,4}$, which does not involve $x_i$, so $j$ imposes no
constraint. For $j\in S_i$ both vertical branches fail, so the pair is legal exactly
when $|x_i-x_j|\ge\bar w_{ij}$, which removes the open interval in
Equation~\eqref{eq:app-slice}. Removing finitely many open intervals from a closed
interval leaves finitely many closed intervals; endpoints are legal because contact
is allowed. The $y_i$ slice follows by exchanging axes. A soft cluster has no
overlap constraint, so its slice is the canvas interval, and a fixed object's slice
is a single point.
\end{proof}
The CD search of Appendix~\ref{app:alg-cd} does not enumerate $I_i$. It keeps the
component that contains the current value, bounded by the nearest blocker on each
side and shrunk by a clearance of $10^{-4}\,\mu$m; a move across a blocker is left to
the proposal stages. Successive moves of other macros change $S_i$, so CD sweeps can
pass between polyhedra $P_\sigma$ even though each single move stays in one interval.

\subsection{Incremental scoring}

Let $\Psi(p)$ be the cache obtained by computing everything from scratch at
placement $p$: pin positions, per-net extrema and, for each net $n$ and object $i$,
the weighted HPWL term $a_n(p)$, the density-load vector $d_i(p)$ of object $i$,
the raw routing vectors $r_n^o(p)$ of net $n$ and the macro-blockage vectors
$b_i^o(p)$ for orientation $o\in\{H,V\}$, with $b_i^o=0$ for soft clusters. The cache
stores the aggregates
\begin{equation}
 T=\sum_na_n(p),\qquad \Lambda=\sum_id_i(p),\qquad R^o=\sum_nr_n^o(p),\qquad M^o=\sum_ib_i^o(p),
 \label{eq:app-cache}
\end{equation}
and every score is a fixed function of these aggregates: normalization, smoothing and
the two tails (Appendix~\ref{app:objective}).

\begin{appremark}[Cache commutation]\label{rem:cache}
Let $U_i$ move object $i$ and $\widehat U_i$ be the cache update that recomputes
$i$'s pins, the extrema of the nets incident to $i$, and then subtracts the old and
adds the new $a_n$, $r_n^o$ for those nets and $d_i$, $b_i^o$ for $i$. In exact
arithmetic, $\widehat U_i\Psi(p)=\Psi(U_ip)$, and the same holds for any sequence of
moves and for reverting them.
\end{appremark}
\begin{proof}
A net term $a_n$ or $r_n^o$ depends only on the positions of the pins of $n$; it
changes under $U_i$ only if $n$ is incident to $i$. The terms $d_j$ and $b_j^o$
depend only on object $j$'s rectangle and change only for $j=i$. Hence, for
instance, $T(U_ip)=T(p)-\sum_{n\ni i}a_n(p)+\sum_{n\ni i}a_n(U_ip)$, which is exactly
the update; the other three aggregates follow in the same way, provided the update
recomputes each incident net's extrema from all its pins when the moved pin was
extreme. Equal aggregates give equal scores. Induction on the number of moves gives
the sequence statement, and a revert applies the same substitution with the roles of
old and new exchanged.
\end{proof}
The remark is conditional on complete dependency tracking and exact arithmetic. In
floating point, subtracting and adding terms need not cancel, and the summation
order differs from a fresh computation; the measured discrepancies are reported in
the main text.

\begin{proof}[Proof of Proposition~\ref{prop:acceptance}]
Let $\widehat J$ be the incremental score and $J$ the canonical score, with
$|\widehat J(S_p)-J(p)|\le\epsilon$ and $|\widehat J(S_q)-J(q)|\le\epsilon$ at the two
endpoints, and suppose $\widehat J(S_q)\le\widehat J(S_p)-\tau$. Adding and
subtracting the incremental values,
\begin{align*}
 J(q)-J(p)
 &=\bigl[J(q)-\widehat J(S_q)\bigr]+\bigl[\widehat J(S_q)-\widehat J(S_p)\bigr]
   +\bigl[\widehat J(S_p)-J(p)\bigr]\\
 &\le\epsilon-\tau+\epsilon=-\tau+2\epsilon .
\end{align*}
If $\tau>2\epsilon$ the right side is negative, so the canonical score strictly
decreases.
\end{proof}
If the three components carry absolute errors at most $\epsilon_W$,
$\epsilon_D$ and $\epsilon_C$, the triangle inequality and
the weights $1,\tfrac12,\tfrac12$ give
$\epsilon=\epsilon_W+\tfrac12\epsilon_D+\tfrac12\epsilon_C$.
The proposition needs the bound at the two states compared; the sampled errors in the
main text satisfy it with a wide margin on the sampled states, which is the sense in
which the CD threshold of $10^{-9}$ and the assignment threshold of $10^{-7}$ are safe.

\subsection{Tail statistics and the cost of a scored move}

For $z\in\mathbb R^N$ and an integer $1\le k\le N$, write $z_{(1)}\ge\dots\ge z_{(N)}$
for the sorted entries and $A_k(z)=k^{-1}\sum_{j=1}^kz_{(j)}$ for the top-$k$ mean;
$(a)_+=\max\{a,0\}$ and
$\mathcal Q_k=\{\beta\in\mathbb R^N:0\le\beta_j\le1/k,\ \sum_j\beta_j=1\}$.

\begin{proof}[Proof of Lemma~\ref{lem:tail}]
\emph{Threshold form.} For any $\theta\in\mathbb R$,
\begin{equation*}
 \sum_{j=1}^kz_{(j)}=k\theta+\sum_{j=1}^k\bigl(z_{(j)}-\theta\bigr)
 \le k\theta+\sum_{j=1}^k\bigl(z_{(j)}-\theta\bigr)_+
 \le k\theta+\sum_{j=1}^N(z_j-\theta)_+ .
\end{equation*}
At $\theta=z_{(k)}$ both inequalities are equalities: the first $k$ differences are
nonnegative, and every entry outside the top $k$ is at most $z_{(k)}$, so its
positive part vanishes. Dividing by $k$ gives
$A_k(z)=\min_\theta\{\theta+k^{-1}\sum_j(z_j-\theta)_+\}$, attained at
$\theta=z_{(k)}$.

\emph{Maximum form.} For $\beta\in\mathcal Q_k$ and $\theta=z_{(k)}$,
\begin{equation*}
 \beta^{\top}z=\theta+\sum_j\beta_j(z_j-\theta)\le\theta+\sum_j\beta_j(z_j-\theta)_+
 \le\theta+\frac1k\sum_j(z_j-\theta)_+=A_k(z),
\end{equation*}
using $\sum_j\beta_j=1$ and $\beta_j\le1/k$. The vector with entries $1/k$ on the
indices of the $k$ largest entries lies in $\mathcal Q_k$ and attains $A_k(z)$, so
$A_k(z)=\max_{\beta\in\mathcal Q_k}\beta^{\top}z$.

\emph{Convexity.} The maximum is attained at one of the finitely many vectors that put
$1/k$ on a $k$-subset, so $A_k$ is a maximum of finitely many linear functions:
convex and piecewise affine, and differentiable wherever the $k$-th and $(k+1)$-th
largest entries differ.

\emph{Stability.} For any $\beta\in\mathcal Q_k$,
$\beta^{\top}(z+\delta)\le A_k(z)+\max_{\beta'\in\mathcal Q_k}\beta'^{\top}\delta$;
maximizing over $\beta$ and exchanging $z$ and $z+\delta$ gives
$|A_k(z+\delta)-A_k(z)|\le\max_{\beta\in\mathcal Q_k}|\beta^{\top}\delta|$. Every
$\beta\in\mathcal Q_k$ has $\lVert\beta\rVert_1=1$, $\lVert\beta\rVert_\infty\le1/k$
and $\lVert\beta\rVert_2^2\le\lVert\beta\rVert_\infty\lVert\beta\rVert_1\le1/k$.
H\"older's
inequality in the three dual pairs $(1,\infty)$, $(2,2)$ and $(\infty,1)$ yields the
bounds $\lVert\delta\rVert_\infty$, $\lVert\delta\rVert_2/\sqrt k$ and
$\lVert\delta\rVert_1/k$.
\end{proof}
Both canonical tails are of this form. For $G\ge10$ the density component is
$\tfrac12A_k(\rho)$ with $N=G$ and $k=\lfloor0.1G\rfloor$: the loads are
nonnegative, so summing the largest $k$ nonzero loads and dividing by $k$ equals
$A_k$ of the full vector with its zeros. The congestion component is $A_{k'}$ of the
$N=2G$ concatenated values with $k'=\lfloor0.05\cdot2G\rfloor$, and $A_1$ (the
maximum) when $k'=0$. Convexity in the
loads does not make the score convex in coordinates, because the loads are
themselves piecewise functions of the coordinates. For the solver the
$\lVert\delta\rVert_1/k$ bound is the useful one: a move changes the loads in few
cells, so its effect on a tail is at most the total load it shifts divided by $k$.

\begin{appremark}[Cost of a scored move]\label{rem:move-cost}
Let $T_{\mathrm{mut}}(i)$ be the work to update object $i$'s pins, its incident nets
and its own grid contributions, and $T_{\mathrm{sel}}(n,k)$ the cost of selecting
the $k$ largest of $n$ values, with $T_{\mathrm{sel}}(0,0)=0$. For $G\ge10$, a call
that moves object $i$ and returns its score costs
\begin{equation*}
 T_{\mathrm{move}}(i)=T_{\mathrm{mut}}(i)+\Theta(G)
 +T_{\mathrm{sel}}\bigl(n_{\mathrm{occ}},\min\{k,n_{\mathrm{occ}}\}\bigr)
 +T_{\mathrm{sel}}(2G,k'),
\end{equation*}
where $n_{\mathrm{occ}}\le G$ counts nonzero density cells. The work is therefore at
least linear in $G$, however few cells the move touches.
\end{appremark}
\begin{proof}
The mutation is local by Remark~\ref{rem:cache}. Producing the score then requires,
from the aggregates of Equation~\eqref{eq:app-cache}: dividing every density load by
the cell area and every routing value by its capacity; smoothing both net-demand maps,
which touches each cell a bounded number of times for the fixed range $s=2$; adding
the blockage maps; and the two selections. Each of the first three steps visits all
$G$ cells, giving the $\Theta(G)$ term. The implementation keeps no incremental
order-statistic structure, so the selections are recomputed in full.
\end{proof}
Full recomputation, by contrast, re-derives every pin position, re-routes every net
and re-rasterizes every object. The microbenchmarks in the main text measure the
saving at the IBM grid sizes ($810\le G\le2{,}255$); they do not establish a cost that
depends on the moved object's degree alone.

\subsection{Negative curvature}

\begin{proof}[Proof of Proposition~\ref{prop:curvature}]
Recall the setting: $f$ is twice continuously differentiable on an open set $U$
containing $p$, $v$ is a unit vector with $v^{\top}\nabla^2f(p)v\le-\eta<0$, the sign
$s\in\{-1,1\}$ satisfies $s\nabla f(p)^{\top}v\le0$, and the segment
$\{p+usv:0\le u\le t\}$ lies in $U$, with the Hessian $L$-Lipschitz on it in the
operator norm. Put
$\phi(u)=f(p+usv)$ for $u\in[0,t]$. Then $\phi'(0)=s\nabla f(p)^{\top}v\le0$ and
$\phi''(u)=v^{\top}\nabla^2f(p+usv)v$, because $s^2=1$. The Lipschitz condition and
$\lVert sv\rVert_2=1$ give $|\phi''(u)-\phi''(0)|\le Lu$, so
$\phi''(u)\le-\eta+Lu$. Taylor's formula with integral remainder gives
\begin{equation*}
 \phi(t)-\phi(0)=t\,\phi'(0)+\int_0^t(t-u)\,\phi''(u)\,du
 \le\int_0^t(t-u)(-\eta+Lu)\,du=-\frac{\eta}{2}t^2+\frac{L}{6}t^3 .
\end{equation*}
For $L>0$ and $0<t\le3\eta/(2L)$, $\tfrac{L}{6}t^3\le\tfrac{\eta}{4}t^2$, so the
decrease is at least $\eta t^2/4$. For $L=0$ the bound is $-\eta t^2/2$ on the whole
segment.
\end{proof}
The bound concerns the unconstrained step. If the favorable sign points out of the
feasible set, it says nothing: $f(x)=x-x^2$ on $[0,\tfrac14]$ has $f'(0)=1$ and
$f''(0)=-2$, yet $x=0$ is its strict constrained minimum because the descent sign
leaves the interval. Lane~B also clips, legalizes and polishes each step before
scoring it with the canonical proxy, so the bound applies to the proposal, not to the
accepted endpoint, and Appendix~\ref{app:curvature} shows that its smoothness
hypothesis fails for the operator Lane~B differentiates.

\begin{appremark}[Eigen-residual bound]\label{rem:residual}
Let $\mathbf H$ be a real symmetric matrix, $v$ a unit vector,
$\hat\lambda=v^{\top}\mathbf Hv$ and $r=\mathbf Hv-\hat\lambda v$. Then
$\operatorname{dist}(\hat\lambda,\operatorname{spec}\mathbf H)\le\lVert r\rVert_2$, and
$\hat\lambda<0$ implies $\lambda_{\min}(\mathbf H)<0$.
\end{appremark}
\begin{proof}
Expand $v=\sum_ja_ju_j$ in an orthonormal eigenbasis, $\mathbf Hu_j=\lambda_ju_j$,
with $\sum_ja_j^2=1$. Then $r=\sum_ja_j(\lambda_j-\hat\lambda)u_j$ and
$\lVert r\rVert_2^2=\sum_ja_j^2(\lambda_j-\hat\lambda)^2\ge\min_j(\lambda_j-\hat\lambda)^2$.
The second claim is the Rayleigh bound $\lambda_{\min}(\mathbf H)\le v^{\top}\mathbf Hv$.
\end{proof}
For the IBM01 probe of Appendix~\ref{app:curvature}, the automatic-differentiation
operator is the Hessian of the branch of $J_B$ active at the endpoint, so it is
symmetric in exact arithmetic, as the remark requires; in float32 it is symmetric
only up to rounding. The residual $3.49\times10^{-7}$ then places an eigenvalue of
that operator within $3.49\times10^{-7}$ of $-0.442458$, and its smallest eigenvalue
is at most $-0.442458$. The remark says nothing about whether the operator is the Hessian
of the scored function, and a small residual can locate an interior eigenvalue rather
than the smallest one.

\subsection{Assignment interactions}

Let $p^0$ be the current placement, $p^{i\leftarrow j}$ the placement with only macro
$i$ moved to slot $j$, and $p^z$ the placement after all moves of an assignment $z$.
For an exact objective $F$, write $\Delta_{ij}=F(p^{i\leftarrow j})-F(p^0)$ and define
the interaction residual
$R_F(z)=F(p^z)-F(p^0)-\sum_{i,j}\Delta_{ij}z_{ij}$, so the joint move improves $F$
exactly when $\sum_{ij}\Delta_{ij}z_{ij}+R_F(z)<0$. For the proxy, the weights give
$R_J=R_W+\tfrac12R_D+\tfrac12R_C$. A net that
contains pins of at most one moved macro contributes nothing to $R_W$,
because its HPWL changes identically under the joint move and under that macro's
single move; shared nets can contribute with either sign. Independent-set matching
in detailed placement restricts each matching to cells that share no nets
\citep{chen2008ntuplace3}, which removes this HPWL residual; Lane~C imposes no such
restriction and draws its macros by individual connectivity, so cluster members can
share nets. Density and congestion interact even across
disjoint nets, through the shared tail selections.

\begin{appremark}[An optimal assignment that worsens HPWL]\label{rem:assignment}
Two movable macros sit at $x_1=2$ and $x_2=12$ on a common horizontal line. Each has
one zero-offset pin on a shared net of weight 2 and a second, coincident pin on a
unit-weight net to a fixed point terminal at its own starting position. Up to the
positive normalization, the HPWL is
\begin{equation*}
 F(x_1,x_2)=2|x_1-x_2|+|x_1-2|+|x_2-12|.
\end{equation*}
With candidate slots 3 and 11, the optimal assignment predicts a change of $-18$ but
realizes $+14$.
\end{appremark}
\begin{proof}
All pins share one vertical coordinate, so only horizontal spans count. Initially
$F(2,12)=2\cdot10+0+0=20$. The four single moves, each with the other macro in place,
give
\begin{equation*}
\begin{aligned}
 F(3,12)&=2\cdot9+1+0=19, &\qquad F(11,12)&=2\cdot1+9+0=11,\\
 F(2,3)&=2\cdot1+0+9=11, &\qquad F(2,11)&=2\cdot9+0+1=19,
\end{aligned}
\end{equation*}
so, with rows for macros 1 and 2 and columns for slots 3 and 11,
\begin{equation*}
 \Delta=\begin{pmatrix}-1&-9\\-9&-1\end{pmatrix}.
\end{equation*}
The two complete assignments have predicted totals $-1-1=-2$ (each macro to its
nearer slot) and $-9-9=-18$ (the swap), so the optimum is the swap,
$(x_1,x_2)=(11,3)$. Its true value is $F(11,3)=2\cdot8+9+9=34$, a realized change of
$34-20=+14$ and a residual $R_F=14-(-18)=32$. The whole residual comes from the shared
net: each single-move entry credits it with a shortening from 10 to~1, while the swap
shortens it only from 10 to~8.

Feasibility: take both macros of width and height $\tfrac12$, centered on the line
$y=\tfrac12$ of a $14\times1$ canvas. Two macros on one horizontal line are legal when
their centers are at least $\tfrac12$ apart. The center gaps are 10 initially, 9, 1,
1 and 9 in the four single moves, and 8 after the swap, so every placement involved
is legal; all centers lie in $[\tfrac14,\tfrac{55}{4}]$, inside the canvas, and the
terminals are points, which impose no overlap constraint. Lane~C's sequential
commit would realize the full swap here, because either order of the two moves is
legal, and its batch check would then reject the batch, since the incremental score
rises.
\end{proof}
Lane~C scores its matrix entries with the incremental evaluator against the current
background, replaces infinite entries before solving, and commits the assigned moves
one at a time, skipping illegal targets (Appendix~\ref{app:alg-lanec}). Its batch
check measures the interactions within the incremental evaluator, and the phase
endpoint is rescored canonically; the example shows why those checks are needed even
when the assignment is solved exactly and every entry is exact.
